

\ifdefined\maindoc
% Do nothing (we are inside main.tex)
\else
\documentclass[12pt]{report}
\usepackage[a4paper,margin=1in]{geometry}
\usepackage[utf8]{inputenc}
\usepackage[T1]{fontenc}
\usepackage{amsmath,amssymb}
\usepackage{booktabs}
\usepackage{setspace}
\usepackage{graphicx}
\usepackage{enumitem}
\usepackage{placeins}
\usepackage{hyperref}
\usepackage{array}
\usepackage[backend=biber,style=ieee,citestyle=numeric-comp,sorting=none]{biblatex}
\addbibresource{../../references.bib}
\setcounter{biburllcpenalty}{7000}
\setcounter{biburlucpenalty}{8000}
\setcounter{biburlnumpenalty}{9000}
\onehalfspacing
\graphicspath{{figures/}}
\begin{document}
\title{Background}
\fi

Infrastructure and process networks are routinely analysed for three separate properties:
whether the system keeps functioning when a component fails, how much it can carry, and when a
dependent sequence of activities finishes. A water distribution network, a power grid, a plant
process line and a construction schedule are frequently the same kind of directed dependency
structure examined three different ways, and the three literatures built around these
properties, reliability computation, capacity and flow analysis, and project scheduling,
developed largely apart from one another, each with its own vocabulary and its own notion of
what counts as an exact answer.

All three properties reduce to a computation over the same kind of object, a directed graph
whose nodes and edges can fail, vary, or take time to traverse. Reliability computation returns
the probability that at least one path still delivers supply from source to demand once every
component is subject to failure. Flow analysis returns the greatest throughput the network
delivers given the maximum capacity each edge and node sustains, and where it bottlenecks.
Scheduling returns the earliest completion time of a dependent sequence of activities and which
of them carry no room to slip. Three different quantities, a probability, a capacity and a
duration, computed over structures the three literatures largely built their own methods for
without citing one another's.

\section{The Methodological Landscape}
\label{sec:bg-landscape}

A reliability block diagram arranges components in series, parallel, and $k$-out-of-$n$ groups
whose system reliability has a closed form \cite{billinton1992network}, and fault and event trees
widen the reachable logic by building the failure mechanism explicitly as a tree of gates
\cite{xing2008fault}, both efficient exactly as far as the failure logic reduces to a tree or a
series-parallel structure. Path and cut-set decomposition reaches arbitrary topologies exactly by
splitting the network into subproblems \cite{liu2012improved,kim2013network}, at a combinatorial
cost that grows with the same redundancy that makes the analysis worth doing, and Monte Carlo
simulation escapes structure entirely in exchange for exactness, at a sample cost that grows
demanding when failures are rare \cite{fishman1986comparison}. Binary decision diagrams compile
the same Boolean structure into a canonical form on which reliability is a single traversal
\cite{bryant1986graph,rauzy1993new}, and cutset conditioning and the junction-tree algorithm reach
the same computation from probabilistic inference instead, treating reachability as marginal
inference on a directed graphical model \cite{pearl1988probabilistic,lauritzen1988local}. Every
one of these methods consumes a point-valued component probability.

Series-parallel reduction collapses independent stretches pair by pair
\cite{satyanarayana1985linear}, factoring conditions on a single component to split the graph
into independent cases \cite{satyanarayana1983factoring}, and polygon-to-chain transformation
widens what can be collapsed in one step beyond a plain series or parallel pair
\cite{wood1985polygon}. Treewidth-bounded tree decomposition generalises the same idea to graphs
no sequence of series-parallel reductions can fully collapse, bounding the cost of exact
computation by a structural width of the graph rather than its raw size
\cite{goharshady2020treewidth}, and the hash-consed representations binary decision diagrams also
depend on store a repeated substructure once rather than rebuilding it
\cite{filliatre2006hashconsing}.

A dependent sequence of activities poses a structurally similar problem under a different
algebra. The critical path method computes the longest path through a project network and the
float every other activity carries, a computation whose float calculus dates to the late 1950s
\cite{kelley1959critical}.

Flow along capacitated edges is the fourth reading, and in its deterministic form it is the
oldest of the four, dating to the max-flow-min-cut results of the 1950s \cite{ford1956}. What
deterministic flow analysis does not address is what happens once the capacities themselves are
uncertain.

\section{Uncertainty in Network Analysis}
\label{sec:bg-uncertainty}

A component's failure probability, an edge's capacity, and an activity's duration are rarely
known exactly. Each is typically estimated from sparse failure data, elicited from expert
judgement, or quoted as a manufacturer's tolerance, and representing that further layer of
not-knowing is itself a choice with several competing answers.

Fuzzy set theory represents such a quantity as a possibility distribution rather than a
probability distribution, assigning each possible value a graded degree of membership rather
than a frequency \cite{zadeh1965fuzzy}. Dempster-Shafer evidence theory assigns belief mass to
sets of possible values directly, combining evidence from independent sources without requiring
a single probability assignment over individual outcomes \cite{shafer1976mathematical}. Imprecise
probability represents the same not-knowing as a bounded set of probability distributions rather
than one distribution or one possibility function: an interval bounds an unknown probability
between two fixed values, and a probability box bounds an entire unknown distribution between two
fixed ones \cite{ferson2003constructing,williamson1990probabilistic}.

Reliability computation's exact methods, from reliability block diagrams through junction trees,
consume a single point-valued probability throughout; the imprecise-probability treatment above
is the direct answer to that limitation.

Flow under uncertainty takes forms that do not coincide with each other. Robust optimisation
treats an uncertain capacity adversarially rather than probabilistically: a capacity is drawn
from an uncertainty set rather than assigned a distribution, and the network is evaluated
against its worst case within that set. The approach has its own complexity landscape, the
robust maximum-flow problem is solvable in polynomial time, while the corresponding robust
minimum-cut problem is NP-hard \cite{bertsimas2013robust}. No probability measure enters this
formulation at any point. The stochastic-flow network reliability problem, founded by Doulliez
and Jamoulle, holds capacities as point-valued random variables instead and asks for the
probability that the network can deliver a stated demand: each arc's capacity is a discrete
random variable, and network reliability comes from decomposing the space of feasible capacity
configurations \cite{doulliez1972transportation}. The method family that follows, state-space
decomposition, boundary-point enumeration, and sum-of-disjoint-products techniques for combining
them, mirrors path and cut-set reliability computation closely enough that the same formalism
travels into project scheduling directly: an activity's duration and cost stand in for an arc's
capacity states, precisely the modelling move behind the exact multi-state project reliability
method discussed in Section~\ref{sec:bg-frameworks} \cite{huang2020exact}.

Duration uncertainty in scheduling followed a similar split. PERT replaced each fixed duration
with a three-point estimate and a distributional assumption almost as soon as the critical path
method itself appeared \cite{malcolm1959pert}. An interval rather than a distribution over each
duration is harder still: deciding whether an activity is possibly critical is NP-complete in
general \cite{chanas2002complexity}.

Across reliability, flow, and schedule, the same handful of representations recur, point-valued
distributions, robust or adversarial bounds, and imprecise probability, but no single
representation is used consistently across all three; each domain settled on its own answer
largely independently of the other two.

\section{Multi-Analysis Methods for Complex Infrastructure}
\label{sec:bg-frameworks}

Reasoning about a complex infrastructure system rarely stops at one question, and running more
than one kind of analysis against a single shared network representation is an established
pattern in engineering software, not unique to any one discipline. \texttt{pandapower}, an
open-source power-system analysis tool, runs power flow, optimal power flow, state estimation,
topological graph search, and short-circuit calculation against one explicit graph of buses,
lines and transformers, a diversity of analysis paradigm matched by commercial power-system
suites such as DIgSILENT PowerFactory and PSS/E \cite{thurner2018pandapower}. The {ArcGIS}
Utility Network models power, gas, and water networks as an explicit graph of features connected
through topology, and runs eight distinct trace types, connectivity, isolation, subnetwork
discovery, loop detection among them, against that one shared topology
\cite{esri_utilitynetwork}. Neither tool's analyses are this thesis's reachability reliability,
capacitated flow, or activity scheduling, and neither propagates interval or probability-box
uncertainty.

Whether an existing tool or method combines reachability reliability, capacitated flow, and
activity scheduling specifically, on one network representation, is a narrower question. The
search for it turned up three distinct partial answers among tools and methods, plus one
comparable piece of doctoral research, and none of the four meets it fully.

Commercial reliability software has, in practice, already paired reliability with flow. Several
of the reliability block diagram and fault tree suites still in active use, ReliaSoft's
BlockSim, Isograph's Availability Workbench, and ITEM Software's ITEM ToolKit among them, bundle
a throughput or process-flow module alongside their core reliability engine, modelling
multi-flow-type capacity, bottlenecks and resource allocation on the same diagrams their
reliability analyses use \cite{reliasoft2024blocksim,isograph2024availability,itemsoftware2024toolkit}.
None of their public documentation describes a project-scheduling or critical-path module, and
all three represent component life and failure with classical distributions, fitted from data or
specified by the analyst, rather than intervals or probability boxes.

Open-source infrastructure simulation platforms have, from the opposite direction, paired flow
with something schedule-shaped. InfraRisk integrates power-flow, water-hydraulic and
traffic-flow simulators for interdependent power, water and road networks under its own banner
of resilience analysis, and its recovery module sequences repair-crew actions with real start
and end times computed from travel time and repair duration, either by heuristic prioritisation
or by model-predictive-control optimisation \cite{balakrishnan2022infrarisk}. That sequencing is
a logistics schedule rather than a precedence-network critical-path computation, since it
optimises the order and timing of repair actions against a performance-loss metric rather than
the float of a fixed activity network. Its reliability side is likewise not a source-to-node
reachability computation. The hazard module derives point component-failure probabilities from
exposure and intensity to drive Monte-Carlo-style disaster scenarios, and the platform reports
its results as simulated performance loss rather than an exact reliability figure.

A third direction pairs reliability with schedule directly, inside project-network reliability
itself. Modelling a project's activities as the arcs of a multi-state network and asking for the
probability that the project completes within both a time and a budget constraint is an
established line of work. Huang, Huang and Lin's decomposition method computes this exact
project reliability, splitting the space of feasible completion-time and budget combinations
into subsets via a project's upper and lower boundary points and evaluating each by
inclusion-exclusion, an approach the authors validate against full enumeration on a real
thirteen-activity manufacturing project \cite{huang2020exact}. The method returns a project
reliability of 0.906 for one tested time-budget pair on that project and traces how the figure
moves as the constraints are relaxed. It has no notion of a physical arc capacity carrying a
transportable flow, activity duration and cost standing in its place instead, and every input is
a discrete state with a fixed probability rather than an interval or a probability box.

Doctoral research comes closer than any commercial or open-source tool to combining two of the
three. Tong's Georgia Tech dissertation develops new methods for infrastructure flow networks
and investigates their reliability in terms of both connectivity and flow capacity together,
using a probability propagation method published the same year and belonging to the same
message-passing lineage surveyed in Section~\ref{sec:bg-landscape} \cite{tong2021infrastructure}.
Schedule and critical-path analysis are absent from it, and the probability treatment throughout
is point-valued rather than interval- or probability-box-based.

No candidate combines all three analyses on one network, and the four above do not even pair the
same two twice: reliability and flow are paired in commercial software with no schedule
dimension, flow and a form of scheduling are paired in open-source simulation platforms with no
formal reliability computation, reliability and schedule are paired in project-network theory
with no true flow network, and reliability and flow are paired again, this time exactly, in
Tong's dissertation, still without schedule. Every one of the four, without exception, represents
its uncertain inputs precisely or adversarially, never as an interval or a probability box, and
neither \texttt{pandapower} nor the {ArcGIS} Utility Network above changes that picture: both run
several kinds of analysis on one network, but neither runs this thesis's specific three, and
neither propagates imprecision either.

\section{Availability of Engineering Analysis Software}
\label{sec:bg-delivery}

In engineering analysis, three patterns of software availability recur, browser-delivered
simulation-as-a-service, licensed desktop installation, and web-based {GIS}, and none of them
has been adopted by the no-code and low-code business-software movement.

The no-code and low-code software category defines itself around business application
development throughout its own literature. A systematic review of 383 definitions across 306
publications from 2014 to 2025 characterises no-code platforms as visual, drag-and-drop
environments that let a non-programmer build software without writing code, and low-code
platforms as the same idea with a professional-developer audience layered on top
\cite{abendroth2026democratization}. Mapping the identified platforms onto ten functional system
classes, the same review names enterprise-resource-planning extensions, customer-relationship
platforms, multi-experience development platforms, business-process automation and management
suites, robotic process automation, document management systems, and integration-platform hubs
as the categories the low-code and no-code movement actually occupies, illustrated with named
commercial examples such as Mendix, Salesforce, UiPath, Bizagi and Zapier. Engineering or
scientific analysis, reliability, flow, structural or otherwise, does not appear among them.

In engineering analysis, browser-delivered simulation-as-a-service is one established
alternative. SimScale, launched in 2013, runs computational fluid dynamics, finite element,
thermal and electromagnetic analysis entirely in the browser with no local installation
\cite{simscale}, though its computation runs in the vendor's cloud rather than on the user's own
machine. The licensed desktop {GUI} reliability suite is a second, older pattern: the same
commercial software named in Section~\ref{sec:bg-frameworks} solves the no-programming half of
the delivery problem as a local install rather than a browser application, the reverse
trade-off from SimScale's.

Some web-based {GIS} platforms exist specifically for hazard and infrastructure risk assessment.
A {PRISMA}-based systematic review of this literature, screening 1{,}775 articles and admitting
65, organises the field around themes including visualisation and user-interface design and
decision-support systems \cite{daud2024webgis}. Repetto et al.'s platform is a concrete member of
it. It delivers wind-monitoring, forecasting and statistical risk mapping for ports, motorways,
rail lines and industrial plants exposed to wind. Its interface is web-based {GIS} built on free
and open geospatial software, run by port authorities and operators who never touch the
underlying numerical wind and wave models themselves \cite{repetto2018windgis}. Most platforms
in this category focus on monitoring and forecasting rather than a structural reliability or
component-failure computation.

Local-first software keeps the usability of the modern web while storing and computing its data
on the user's own machine rather than a remote server \cite{kleppmann2019local}, an architecture
that matters most in exactly the industries whose proprietary system models cannot leave the
premises.

SimScale is browser-based but not local; the desktop suites are local but not browser-based; the
web-{GIS} platforms are browser-based but run on a remote server. No tool found in this search
combines browser delivery, no programming, and local computation and data all at once.

\section{Julia as a Language for Network Modelling and Analysis}
\label{sec:bg-language}

{MATLAB} has offered built-in graph and network functions since the mid-2010s: shortest-path,
maximum-flow, centrality, connected-component and spanning-tree computations all operate on its
native \texttt{graph} and \texttt{digraph} classes \cite{mathworks_graphalgorithms}. R's
equivalent general-purpose graph package is igraph, used for shortest paths, centrality,
community detection and flow computation over one shared graph object \cite{csardi2006igraph},
and R additionally carries a package purpose-built for structural reliability specifically,
computing system and survival signatures over a network's topology \cite{aslett_reliabilitytheory},
the kind of dedicated reliability tooling neither {MATLAB}'s built-in functions nor Python's
ecosystem, below, provide natively. Rust's graph ecosystem is newer and smaller, centred on
\texttt{petgraph}, a general-purpose graph data-structure library with the same shortest-path,
spanning-tree and traversal algorithms as the languages above \cite{petgraph}.

Python's equivalent general-purpose graph package is NetworkX, whose own canonical paper
presents it as covering simple graphs, directed graphs, and graphs with parallel edges, over a
single native numeric type \cite{hagberg2008networkx}. No single Python package purpose-built for
network reliability computation, in the reliability block diagram, path-set, or decision-diagram
sense this literature uses the term, was found in this search, the way R's
\texttt{ReliabilityTheory} package is; the closest candidate, \texttt{pgmpy}, is a general
Bayesian-network and probabilistic-graphical-model toolkit with belief propagation and structure
learning built in, useful machinery for the same class of problem but not a package built to
solve it directly. This is a gap in packaged tooling specifically, not in published research:
network reliability analysis is itself an actively reviewed field with applications across
communications, transportation, and electrical and gas distribution \cite{gaur2021literature},
and that research draws on whichever language a given group already works in without converging
on one shared package the way NetworkX did for general graph algorithms.

Julia's own graph ecosystem, centred on \texttt{Graphs.jl} and its satellite packages for
weighted graphs, property graphs and bindings to established external libraries, occupies the
equivalent general-purpose position to NetworkX, igraph and \texttt{petgraph}, though no single
paper describes that ecosystem the way Hagberg, Schult and Swart's paper describes NetworkX.
None of the packages named above, in any of the five languages, carries a native notion of an
interval- or probability-box-valued edge or node, which is not a limitation any of them sets out
to address, since none is built for propagating uncertain quantities across a network in the
first place.

Multiple dispatch, choosing which method of a generic function to run from the runtime types of
every argument rather than just the first, is presented in the canonical Julia paper as the
mechanism that lets one generic function definition serve arbitrarily many numeric and
non-numeric types without its author anticipating any of them in advance
\cite{bezanson2017julia}. The same paper grounds Julia's central design claim, that a dynamic
language need not trade its ease of use for the performance a statically compiled one gets, in
the same type system and dispatch mechanism. This is the mechanism this framework's own
propagation step leans on to specialise correctly whether the value flowing through it is a
float, an interval, or a probability box, without the author of the propagation step anticipating
any of the three in advance.

\section{Summary}
\label{sec:bg-summary}

Reliability, flow, and schedule analysis grew as three separate literatures, each with its own
exact methods and its own way of representing what it does not know precisely about its inputs.
Running several kinds of analysis against one shared network object is not itself a new idea in
engineering software, but no tool, method, or doctoral thesis found in this search combines
reachability reliability, capacitated flow, and activity scheduling on one network
representation, with any of them propagating imprecision natively across the combination. No
engineering-analysis tool found in this search delivers that combination through an interface
that asks its user for files and clicks rather than code, either. This thesis's own framework is
built to close that gap. It computes all three analyses from one network model, propagates
interval uncertainty natively through each and probability-box uncertainty through reliability
specifically, and reaches its user through a local, no-code interface rather than a script.

\ifdefined\maindoc
% do nothing
\else
\printbibliography
\end{document}
\fi
